\documentclass[11pt,reqno]{amsart}
\usepackage[T1]{fontenc}
\usepackage[utf8]{inputenc}
\usepackage{lmodern}
\usepackage{amsmath,amssymb,amsfonts,mathrsfs,mathtools,amsthm}
\usepackage{geometry,enumitem,booktabs,longtable,array,microtype}
\usepackage[dvipsnames]{xcolor}
\usepackage[pagebackref=true]{hyperref}
\usepackage{bookmark,pdfcomment}
\usepackage{listings}
\geometry{margin=1in}
\numberwithin{equation}{section}
\setcounter{tocdepth}{3}
\setlength{\emergencystretch}{2em}
\hypersetup{colorlinks=false,pdfborder={0 0 1},linkbordercolor={1 0 0},citebordercolor={0 0.3 1},urlbordercolor={0 0.6 0.6},linktoc=all,bookmarksnumbered=true,bookmarksopen=true,bookmarksopenlevel=1,pdftitle={Prime-Corrected Clocks and Synthetic Zeta Functions: Local Rouche Certificates and Rigorous Indexing},pdfauthor={Rotkiv Aynasac}}
\makeatletter
\let\CG@origref\ref
\let\CG@origeqref\eqref
\newcommand{\CG@gettip}[1]{\@ifundefined{CGTip@#1}{Reference: #1}{\csname CGTip@#1\endcsname}}
\DeclareRobustCommand{\CGref}[1]{\pdftooltip{\CG@origref{#1}}{\CG@gettip{#1}}}
\DeclareRobustCommand{\CGeqref}[1]{\pdftooltip{\CG@origeqref{#1}}{\CG@gettip{#1}}}
\expandafter\gdef\csname CGTip@thm:adjacent\endcsname{If every adjacent pair has combined center error below the synthetic gap, ordinates are all distinct and RH plus simplicity follows.}
\expandafter\gdef\csname CGTip@thm:cell\endcsname{The directional Voronoi inequality at every index implies RH plus simplicity.}
\expandafter\gdef\csname CGTip@prop:symmetry\endcsname{Synthetic zeta has exact functional equation and a pole residue fixed by the entire synthetic product.}
\expandafter\gdef\csname CGTip@prop:loewner\endcsname{The synthetic Stieltjes transform has a strictly positive Loewner divided-difference kernel.}
\expandafter\gdef\csname CGTip@thm:flow\endcsname{The prime-only analytic scalar flow preserves strict order for 0 < a <= 2.}
\expandafter\gdef\csname CGTip@thm:rational\endcsname{Exact rational frequencies with Moebius coefficients b_X(n/d)=mu(d)c_X(d)/sqrt(nd).}
\expandafter\gdef\csname CGTip@thm:two\endcsname{Two small Rouché discs and one large disk certify exactly two real simple roots under two complex remainder bounds.}
\expandafter\gdef\csname CGTip@thm:classification\endcsname{A local critical point is classified as a real pair, conjugate nonreal pair, or an exact double zero under a validated complex remainder bound.}
\expandafter\gdef\csname CGTip@thm:one\endcsname{A validated Taylor inequality with a linear comparator proves one real simple zero in a disk.}
\expandafter\gdef\csname CGTip@prop:index\endcsname{The Arb or Turing-indexed ball must lie inside the same disk certified by Rouche, especially after Newton recentering.}
\expandafter\gdef\csname CGTip@thm:hankel\endcsname{Hermite-Hankel positivity on a conjugation symmetric contour characterizes reality, and definite positivity characterizes simple real roots.}
\expandafter\gdef\csname CGTip@thm:isolated\endcsname{Fourth-order Taylor roots are accurate relative to gap under a strong analytic factor isolation hypothesis.}
\expandafter\gdef\csname CGTip@prop:portmanteau\endcsname{Weak spacing limits control density but cannot exclude infinitely many exceptional coincident ordinates.}
\expandafter\gdef\csname CGTip@eq:I\endcsname{First strict Rouché inequality on the two small discs.}
\expandafter\gdef\csname CGTip@eq:II\endcsname{Second strict Rouché inequality on the enclosing large disc.}
\expandafter\gdef\csname CGTip@eq:oneineq\endcsname{One-root Rouche inequality involving Taylor coefficients and a whole-disc remainder bound.}
\expandafter\gdef\csname CGTip@eq:samedisc\endcsname{Same-disc ball inclusion required for rigorous matching to the indexed zeta zero.}
\expandafter\gdef\csname CGTip@eq:adjacent\endcsname{Each adjacent pair total absolute ordinate error is below the synthetic gap.}
\expandafter\gdef\csname CGTip@eq:fourier\endcsname{The rational-frequency expansion of the prime Euler phase.}
\makeatother
\newtheorem{theorem}{Theorem}[section]
\newtheorem{proposition}[theorem]{Proposition}
\newtheorem{lemma}[theorem]{Lemma}
\newtheorem{corollary}[theorem]{Corollary}
\theoremstyle{definition}
\newtheorem{definition}[theorem]{Definition}
\newtheorem{hypothesis}[theorem]{Hypothesis}
\theoremstyle{remark}
\newtheorem{remark}[theorem]{Remark}
\newcommand{\R}{\mathbb R}
\newcommand{\C}{\mathbb C}
\newcommand{\dd}{\,\mathrm d}
\newcommand{\Res}{\operatorname{Res}}
\newcommand{\sgn}{\operatorname{sgn}}
\newcommand{\supp}{\operatorname{supp}}
\newcommand{\eps}{\varepsilon}
\newcommand{\XiC}{\Xi_c}
\newcommand{\Var}{\operatorname{Var}}
\lstset{basicstyle=\ttfamily\scriptsize,breaklines=true,breakatwhitespace=true,columns=fullflexible,frame=single}
\begin{document}
\title[Prime-corrected clocks and local certificates]{Prime-Corrected Clocks, Synthetic Zeta Functions, and Certified Local Zero Localization: Arithmetic Identities, Rouch\'e Tests, and Rigorous Indexing}
\author{Rotkiv Ayn\"asac}
\date{23 October 2026}
\begin{abstract}
We develop a prime-only construction of increasing reference ordinates, a corresponding synthetic zeta function, and exact arithmetic identities for its finite-Euler phase. We distinguish this synthetic construction from rigorous local verification of the completed Riemann zeta function. A one-zero Rouch\'e theorem is proved with a validated Taylor remainder and a separate ordinal-index condition based on an independently certified Turing count. In particular, if a Taylor disc is recentered by a Newton step, the Turing-indexed zero must be contained in that \emph{same} disc; containment in an inflated disc about the old center alone is insufficient. We retain a two-zero classification certificate for real and nonreal conjugate pairs, a direct per-index ordinate separation theorem, and the exact rational-frequency M\"obius decomposition of the finite Euler phase. Two previously reported local pair examples near heights 1283 and 1290 are checked independently in ordinary high precision, but their third-derivative bounds have not been independently recertified by Arb/FLINT. All statements about the reality of zeros are local or conditional; neither universal coverage nor the Riemann hypothesis is established.
\end{abstract}
\maketitle
\tableofcontents

\section{Analytic and arithmetic setup}\label{sec:setup}
Write the nontrivial zeros of the Riemann zeta function as
\begin{equation}\label{eq:zeros}
 \rho_k=\beta_k+i\gamma_k,\qquad 0<\beta_k<1,\qquad
 0<\gamma_1\le\gamma_2\le\cdots,
\end{equation}
with \emph{all} positive ordinates included and zeros counted with multiplicity. In particular, no restriction $\beta_k=\tfrac12$ is implicit in this notation. On the real line let
\begin{equation}\label{eq:theta}
 \vartheta(t)=\Im\log\Gamma\left(\frac14+\frac{it}{2}\right)-\frac t2\log\pi,
 \qquad N_0(t)=1+\frac{\vartheta(t)}\pi.
\end{equation}
Here and below logarithmic branches are chosen continuously in the regions at issue. When $N_0$ is strictly increasing, its half-integer levels define the smooth clock
\begin{equation}\label{eq:tau}
 \tau_k=N_0^{-1}\left(k-\frac12\right).
\end{equation}
A useful elementary asymptotic surrogate is
\begin{equation}\label{eq:surrogate}
 \mathcal N(t)=\frac{t}{2\pi}\left(\log\frac{t}{2\pi}-1\right)+\frac78,
 \qquad L(t)=\mathcal N'(t)=\frac{\log(t/2\pi)}{2\pi}.
\end{equation}
The two clocks are not identical. The finite Euler phase is
\begin{equation}\label{eq:phase}
 P_X(s)=\prod_{p\le X}(1-p^{-s})^{-1},\quad
 S_X(t)=-\frac1\pi\sum_{p\le X}\arg(1-p^{-1/2-it}),\quad
 F_X(t)=N_0(t)+S_X(t).
\end{equation}
For each fixed $X$ these functions are computable from the prime list and the index, without the genuine ordinates $\gamma_k$. The half-integer equation $F_X(t)=k-1/2$ need not have a unique solution at an arbitrary large cutoff: a branch-selection rule or a certified continuation is necessary.

A \emph{synthetic system} is a strictly increasing positive sequence $c_1<c_2<\cdots$ tending to infinity with $\sum_kc_k^{-2}<\infty$. We allow any prime-only construction satisfying this definition, including the Abel prime flow of Section~\CGref{sec:flow}. The file \texttt{zCasanya.bin} denotes the separate finite-cutoff data construction with cutoffs $(8,8,512)$, \emph{not} the Abel flow and not a theorem of exact ordinate recovery.

\section{A per-index separation theorem}\label{sec:separation}
Put $e_k=c_k-\gamma_k$ and $d_k=c_{k+1}-c_k>0$. The next theorem exploits the functional-equation symmetry of $\zeta$ without assuming RH.
\begin{theorem}[Adjacent error and spacing criterion]\label{thm:adjacent}
If for every $k\ge1$,
\begin{equation}\label{eq:adjacent}
 \left|c_k-\gamma_k\right|+\left|c_{k+1}-\gamma_{k+1}\right|<c_{k+1}-c_k,
\end{equation}
then every nontrivial zero of $\zeta$ is on the critical line, and every such zero is simple.
\end{theorem}
\begin{proof}
Writing $\Delta_k=\gamma_k-c_k$, we have the exact relation
\begin{equation}\label{eq:gapidentity}
 \gamma_{k+1}-\gamma_k=d_k+\Delta_{k+1}-\Delta_k
 \ge d_k-|\Delta_k|-|\Delta_{k+1}|>0.
\end{equation}
Thus no two zeros in the full multiplicity-counting list have the same positive ordinate. If a zero $\beta+i\gamma$ had $\beta\ne1/2$, the symmetry $\rho\mapsto1-\overline\rho$ would give another zero $1-\beta+i\gamma$ with the same ordinate. A multiple critical-line zero would likewise contribute a repeated ordinate. Both contradict \CGeqref{eq:gapidentity}.
\end{proof}
\begin{theorem}[Directional cell criterion]\label{thm:cell}
For $k\ge2$ put $\ell_k=c_k-c_{k-1}$ and $r_k=c_{k+1}-c_k$. The ordinate $\gamma_k$ belongs to the open Voronoi cell of $c_k$ exactly when
\begin{equation}\label{eq:directional}
 -\frac{r_k}{2}<e_k<\frac{\ell_k}{2}.
\end{equation}
If this is true for all indices, with the first cell treated separately at the lower endpoint, then RH and simplicity follow.
\end{theorem}
\begin{proof}
The cell is $((c_{k-1}+c_k)/2,(c_k+c_{k+1})/2)$. Rearrangement gives the displayed inequalities. The open cells are disjoint; hence the ordinates are distinct and Theorem~\CGref{thm:adjacent}'s final symmetry argument applies.
\end{proof}
\begin{remark}[Non-equivalence of sufficient conditions]\label{rem:criteria}
Theorems~\CGref{thm:adjacent} and~\CGref{thm:cell} give distinct sufficient conditions. A failure of the own-cell condition at one center need not contradict the adjacent-error condition. Neither is established for every index by known mean-square estimates. One repeated-ordinate pair contributes only $d_k^2/(2N)$ to an $N$-point mean-square error if placed at the midpoint between its synthetic neighbors, so an averaged bound cannot exclude such an exceptional pair.
\end{remark}

\section{Synthetic zeta, functional equation, and positive kernels}\label{sec:synthetic}
For a synthetic system define the even entire canonical product and the completed-to-zeta conversion factor
\begin{equation}\label{eq:synthetic}
 \Xi_c(z)=\prod_{k\ge1}\left(1-\frac{z^2}{c_k^2}\right),\quad
 D(s)=\frac12s(s-1)\pi^{-s/2}\Gamma(s/2),\quad
 \zeta_c(s)=\frac{\Xi_c(-i(s-1/2))}{D(s)}.
\end{equation}
The product converges normally on compact subsets and has simple real zeros at $\pm c_k$.
\begin{proposition}[Exact synthetic symmetry]\label{prop:symmetry}
The meromorphic function $\zeta_c$ satisfies
\begin{equation}\label{eq:func}
 \zeta_c(s)=\chi(s)\zeta_c(1-s),\qquad
 \chi(s)=\pi^{s-1/2}\frac{\Gamma((1-s)/2)}{\Gamma(s/2)}.
\end{equation}
Its residue at $s=1$ is $2\Xi_c(i/2)$, so the residue-one normalization is
\begin{equation}\label{eq:norm}
 \zeta_c^{\mathrm{nat}}(s)=\frac{\zeta_c(s)}{2\Xi_c(i/2)}.
\end{equation}
The synthetic functional equation does not imply $\zeta_c=\zeta$.
\end{proposition}
\begin{proof}
Evenness of $\Xi_c$ yields $\Xi_c(-i(s-1/2))=\Xi_c(-i((1-s)-1/2))$. Direct cancellation of Gamma factors gives $D(1-s)/D(s)=\chi(s)$. Since $D(s)/(s-1)\to1/2$ at $s=1$, the residue is $2\Xi_c(-i/2)=2\Xi_c(i/2)$.
\end{proof}
\begin{proposition}[Loewner--Pick kernel of the synthetic product]\label{prop:loewner}
For $w>0$ put $D_c(w)=\prod_k(1+w/c_k^2)$ and
\begin{equation}\label{eq:Stieltjes}
 G_c(w)=-\frac{D_c'(w)}{D_c(w)}=-\sum_k\frac1{w+c_k^2}.
\end{equation}
For distinct positive $x_1,\ldots,x_n$, the Loewner matrix with diagonal values given by derivatives obeys
\begin{equation}\label{eq:Loewner}
 \sum_{i,j=1}^n v_iv_j\frac{G_c(x_i)-G_c(x_j)}{x_i-x_j}
 =\sum_{k\ge1}\left(\sum_{i=1}^n\frac{v_i}{x_i+c_k^2}\right)^2>0
\end{equation}
for any nonzero real vector $v$, with the divided difference at $i=j$ interpreted as $G_c'(x_i)$.
\end{proposition}
\begin{proof}
Subtract two simple fractions in \CGeqref{eq:Stieltjes}. If the sum on the right vanished, the rational function $\sum_i v_i/(w+x_i)$ would vanish at infinitely many distinct $w=c_k^2$ and hence identically; its simple-pole residues $v_i$ must all vanish.
\end{proof}
The kernel holds for arbitrary synthetic increasing sequences and does not compare their zeros with the true $\zeta$ zeros. A further, as yet unproved, transfer inequality would be needed to obtain information about RH.

\section{Analytic prime correction and ordered flow}\label{sec:flow}
For fixed $0<a\le2$ define the Abel-smoothed prime correction
\begin{equation}\label{eq:abel}
 Q_a(t)=\left(\frac t{2\pi}\right)^a,\qquad
 \mathcal S_a(t)=-\frac1\pi\sum_p\sum_{r\ge1}
 \frac{e^{-p^r/Q_a(t)}}{r p^{r/2}}\sin(rt\log p).
\end{equation}
All primes and prime powers occur, but the series and its real derivatives converge locally normally for $t>0$.
\begin{proposition}[Analyticity and correction size]\label{prop:abel}
For fixed $a>0$, $\mathcal S_a$ is real analytic for $t>0$ and
\begin{equation}\label{eq:abelbound}
 |\mathcal S_a(t)|\ll_a\frac{t^{a/2}}{\log t},\qquad
 \left|\tau_k-\frac{\mathcal S_a(\tau_k)}{L(\tau_k)}-\tau_k\right|
 \ll_a\frac{\tau_k^{a/2}}{(\log\tau_k)^2}.
\end{equation}
These are bounds relative to the RvM clock, not to the true ordinates.
\end{proposition}
\begin{proof}
The exponential factor gives local normal convergence even after differentiating. Bounding the $r=1$ prime sum absolutely and using $\pi(u)\ll u/\log u$ gives $\sum_p e^{-p/Q}/\sqrt p\ll\sqrt Q/\log Q$ as $Q\to\infty$. Prime powers $r\ge2$ contribute $O(1+\log\log Q)$, which is smaller. Divide by $L(t)\asymp\log t$ for the second inequality.
\end{proof}
The raw Newton correction need not preserve order. Instead let
\begin{equation}\label{eq:flow}
 V_a(t)=-\frac{(1-2\pi/t)^2}{L(t)}\mathcal S_a(t),\quad
 \partial_\lambda Y_a(\lambda,\nu)=V_a(Y_a(\lambda,\nu)),\quad
 Y_a(0,\nu)=\tau(\nu),
\end{equation}
where $\tau$ is chosen as an increasing high-half-line clock and its initial values exceed $2\pi$.
\begin{theorem}[Global ordered prime flow]\label{thm:flow}
For every $0<a\le2$, equation \CGeqref{eq:flow} has a unique positive-half-line solution for every finite real $\lambda$, and
\begin{equation}\label{eq:flowderivative}
 \partial_\nu Y_a(1,\nu)=\tau'(\nu)
 \exp\left(\int_0^1 V_a'(Y_a(v,\nu))\,\dd v\right)>0.
\end{equation}
Thus the centers $c_{a,k}=Y_a(1,k)$ are strictly ordered and form a synthetic system.
\end{theorem}
\begin{proof}
The double zero of $(1-2\pi/t)^2$ cancels the simple zero of $L(t)$ at $2\pi$ and leaves a locally Lipschitz vector field with equilibrium at $2\pi$. A solution starting above $2\pi$ cannot cross this equilibrium by uniqueness. At infinity, $|V_a(t)|\ll_a t^{a/2}/\log^2t=O(t)$, ruling out escape in finite flow time. Distinct scalar solution curves do not intersect. The displayed derivative is the standard variational formula. Since $a/2<1$ for $a<2$ and the $a=2$ estimate is $O(t/\log^2t)$, the displacement is $o(t)$ at infinity. As $\tau_k\asymp k/\log k$, the resulting centers satisfy $c_k\asymp k/\log k$ and $\sum c_k^{-2}<\infty$.
\end{proof}
\begin{remark}[Exponent not optimized]\label{rem:alpha}
One convenient analytic parameter is $a=1/\sqrt2$, but minimizing the auxiliary expression $\log(1/a)+a^2$ has no demonstrated implication for minimizing $|c_{a,k}-\gamma_k|$. This paper does not claim an optimal exponent.
\end{remark}

\section{Exact arithmetic phase and rational resonances}\label{sec:rational}
Assume for this section that a fixed-cutoff clock $F_X$ is strictly increasing on a specified interval $(T,2T)$, and write $\kappa_k=F_X^{-1}(k-1/2)$ for its levels in that interval. The Poisson identity is an identity of distributions:
\begin{equation}\label{eq:poisson}
 \sum_{\kappa_k\in(T,2T)}\delta_{\kappa_k}(t)
 =F_X'(t)\left[1+2\sum_{m\ge1}(-1)^m\cos(2\pi mF_X(t))\right].
\end{equation}
The equality is tested against compactly supported smooth functions, for which the Fourier series on the right converges after integration. For $w\in C_c^\infty((T,2T))$ with even extension $w_e(t)=w(|t|)$,
\begin{equation}\label{eq:trace}
 \sum_k w(\kappa_k)=\int_T^{2T}w(t)N_0'(t)\dd t
 -\frac1{2\pi}\sum_{p\le X}\sum_{r\ge1}
 \frac{\log p}{p^{r/2}}\widehat w_e(r\log p)+\mathcal R_X[w],
\end{equation}
where
\begin{equation}\label{eq:remainder}
 \mathcal R_X[w]=2\sum_{m\ge1}(-1)^m\int_T^{2T}w(t)F_X'(t)\cos(2\pi m F_X(t))\dd t.
\end{equation}
The Fourier coefficients of $w\circ F_X^{-1}$ decay faster than any power of $m$, proving convergence of \CGeqref{eq:remainder}. None of these formulas asserts that $\kappa_k$ equals a genuine zero ordinate.

Define $P(X)=\prod_{p\le X}p$, $d\mid P(X)$ squarefree, $n$ supported on primes $\le X$, $(n,d)=1$, and
\begin{equation}\label{eq:weights}
 c_X(d)=\prod_{\substack{p\le X\\p\nmid d}}\left(1-\frac1p\right),\qquad
 b_X(n/d)=\frac{\mu(d)c_X(d)}{\sqrt{nd}}.
\end{equation}
\begin{theorem}[M\"obius-weighted rational Fourier expansion]\label{thm:rational}
For fixed finite $X$, with absolute convergence,
\begin{equation}\label{eq:fourier}
 e^{-2\pi i S_X(t)}=\sum_{\substack{q=n/d>0\\d\mid P(X),\ (n,d)=1}}b_X(q)q^{it},
 \qquad
 \sum_q|b_X(q)|=\prod_{p\le X}(1+2p^{-1/2}).
\end{equation}
The denominator-frequency coefficient is $b_X(1/d)=\mu(d)c_X(d)/\sqrt d$; only after a stationary-phase factor $\sqrt q$ does it become $b_X(1/d)\sqrt{1/d}=\mu(d)c_X(d)/d$.
\end{theorem}
\begin{proof}
For $u=p^{-1/2}$ and $w=p^{it}$,
\begin{equation}\label{eq:localfactor}
 \frac{1-uw^{-1}}{1-uw}=-uw^{-1}+(1-u^2)\sum_{r\ge0}u^rw^r.
\end{equation}
The Euler phase is the product of these ratios over $p\le X$. Multiplying and reducing $n/d$ leaves $\mu(d)c_X(d)/\sqrt{nd}$, and each local absolute coefficient sum is $1+2u$. The explicit stationary-phase weight is obtained by multiplication by $\sqrt q$; confusing the two would change the theorem.
\end{proof}
The arithmetic sum \CGeqref{eq:fourier} presents a possible target for a nontrivial cancellation theorem uniform in the growing prime cutoff. Its absolute coefficient mass is insufficient to conclude such cancellation.

\section{The completed zeta and a two-zero classification}\label{sec:pair}
Define the real entire function
\begin{equation}\label{eq:completed}
 f(z)=\xi\left(\frac12+iz\right),\qquad
 \xi(s)=\frac12s(s-1)\pi^{-s/2}\Gamma(s/2)\zeta(s).
\end{equation}
It is even and satisfies $f(\overline z)=\overline{f(z)}$. A zero $\rho=\beta+i\gamma$ becomes a zero
\begin{equation}\label{eq:coordinates}
 z=\gamma+i\left(\frac12-\beta\right).
\end{equation}
Thus off-critical zeros yield nonreal conjugate zeros of $f$ with the same real coordinate $\gamma$. The use of $\xi$ in certification is explicit: the synthetic centers may be constructed without it, but a finite validation test examines the true function.

\begin{theorem}[Two-disc Rouch\'e certificate]\label{thm:two}
Let $f$ be holomorphic on a neighborhood of $\overline D(b,3h/2)$ and real on the real axis, with $b\in\R$, $h>0$, $f''(b)\ne0$. Put
\begin{equation}\label{eq:quad}
 Q(z)=\frac{f''(b)}2\left((z-b)^2-h^2\right),\qquad
 M_3=\sup_{|z-b|\le3h/2}|f'''(z)|.
\end{equation}
If the two strict inequalities
\begin{align}
 \left|f(b)+\tfrac12f''(b)h^2\right|+\tfrac54h|f'(b)|
 +\tfrac{125}{384}M_3h^3&<\tfrac7{32}|f''(b)|h^2,\label{eq:I}\\
 \left|f(b)+\tfrac12f''(b)h^2\right|+\tfrac32h|f'(b)|
 +\tfrac9{16}M_3h^3&<\tfrac58|f''(b)|h^2\label{eq:II}
\end{align}
are satisfied, then $f$ has one real simple zero in each of $D(b-h,h/4)$ and $D(b+h,h/4)$, and exactly two zeros, counted with multiplicity, in $D(b,3h/2)$.
\end{theorem}
\begin{proof}
Taylor's theorem on a complex disk gives
\begin{equation}\label{eq:quadrest}
 |f(z)-Q(z)|\le\left|f(b)+\tfrac12f''(b)h^2\right|+|f'(b)|\,|z-b|
 +\tfrac{M_3}{6}|z-b|^3.
\end{equation}
On $|z-(b\pm h)|=h/4$ one has $|Q(z)|\ge 7|f''(b)|h^2/32$ and $|z-b|\le5h/4$; \CGeqref{eq:I} establishes the strict Rouch\'e bound. On $|z-b|=3h/2$, the corresponding estimates are $|Q(z)|\ge5|f''(b)|h^2/8$ and \CGeqref{eq:II}. Each small circle contains one zero and the large circle exactly two. Both small circles are conjugation invariant, so uniqueness forces their roots to be real; multiplicity one follows from their zero count.
\end{proof}
\begin{theorem}[Classification at a nondegenerate critical point]\label{thm:classification}
Let $f'(b)=0$, $f''(b)\ne0$ and $f(b)\ne0$; set
\begin{equation}\label{eq:hcrit}
 h=\sqrt{\frac{2|f(b)|}{|f''(b)|}},\qquad
 \sigma=\sgn(f(b)f''(b)).
\end{equation}
If
\begin{equation}\label{eq:84}
 \frac{hM_3}{|f''(b)|}<\frac{84}{125},\qquad
 M_3=\sup_{|z-b|\le3h/2}|f'''(z)|,
\end{equation}
then $D(b,3h/2)$ contains precisely two simple zeros: real, located one each in $D(b\pm h,h/4)$ if $\sigma=-1$, or nonreal conjugates located one each in $D(b\pm ih,h/4)$ if $\sigma=+1$. If $f(b)=f'(b)=0\ne f''(b)$ and $r\sup_{|z-b|\le r}|f'''(z)|<3|f''(b)|$, then $b$ is the only zero in $D(b,r)$ and has multiplicity exactly two.
\end{theorem}
\begin{proof}
The quadratic Taylor polynomial is $f''(b)((z-b)^2+\sigma h^2)/2$. The small-circle comparison gives the numerical threshold $84/125$, and the large-circle comparison needs only $10/9$. Hence Rouch\'e proves the two counts. For $\sigma=-1$ the invariant small disks contain real roots; for $\sigma=+1$ they are off-axis disjoint disks exchanged by conjugation, so the roots are a nonreal pair. In the double case, comparison with $f''(b)(z-b)^2/2$ establishes a total multiplicity two at the already-known double zero.
\end{proof}
The sign test is \emph{not} valid in isolation: the strict complex-analytic remainder bound is essential. These local results do not exclude additional nonreal zeros outside the certified disks.

\section{Rouch\'e's one-zero certificate and Turing indexing}\label{sec:one}
In large-scale local verification it can be cheaper to certify one zero at a time. For numerical stability one may work with the holomorphic continuation of Hardy's $Z$ function
\begin{equation}\label{eq:hardy}
 Z(z)=e^{i\vartheta(z)}\zeta\left(\tfrac12+iz\right),
\end{equation}
where the branch of $\vartheta(z)$ is analytically continued from the real line on a disk avoiding the poles and branch singularities of the Gamma terms. Such a continuation is real on the real axis and has the same local zero set, with multiplicity, as $f$ in \CGeqref{eq:completed}. The next theorem applies to either $f$, $Z$, or a nonvanishing analytic rescaling.
\begin{theorem}[Validated Taylor disk containing one real simple zero]\label{thm:one}
Let $g$ be holomorphic on a neighborhood of $\overline D(b,r)$ with $b\in\R$, and real on the real axis. Fix an integer $K\ge2$ and let $a_j=g^{(j)}(b)/j!$. Define
\begin{equation}\label{eq:onebound}
 M_K=\frac1{K!}\sup_{|z-b|\le r}|g^{(K)}(z)|,
 \quad S_K(r)=\sum_{j=2}^{K-1}|a_j|r^j+M_Kr^K.
\end{equation}
If $a_1\ne0$ and
\begin{equation}\label{eq:oneineq}
 \boxed{|a_0|+S_K(r)<|a_1|r},
\end{equation}
then $g$ has exactly one zero in $D(b,r)$, counted with multiplicity, and it is real and simple.
\end{theorem}
\begin{proof}
Taylor's formula along a complex line segment gives $|g(z)-(a_0+a_1(z-b))|\le S_K(r)$ on $|z-b|=r$. By the reverse triangle inequality, $|a_1(z-b)|-|a_0|\ge |a_1|r-|a_0|>S_K(r)$. Hence Rouch\'e compares $g$ with $a_1(z-b)$, which has exactly one simple zero in $D(b,r)$. Reality follows because a nonreal zero would entail a distinct conjugate zero in the same disk.
\end{proof}
\begin{proposition}[Correct ordinal-index transfer]\label{prop:index}
Suppose Theorem~\CGref{thm:one} validates the disk $D(b,r)$ and an independent, rigorously indexed Turing/Arb routine returns a complex ball $\mathcal B_k$ containing the $k$th zeta zero $1/2+i\gamma_k$ and certifies that $\gamma_k$ is its true ordinal index (with the relevant multiplicity convention). If
\begin{equation}\label{eq:samedisc}
 \boxed{\mathcal B_k\subset\{s:\ |s-(1/2+ib)|<r\}},
\end{equation}
then the single root certified by Rouch\'e is exactly the root carrying index $k$. If the disc was recentered at a real Newton point $b_1$, the containment test \emph{must use $D(b_1,r_1)$}, not the larger disk $D(x,r_1+|b_1-x|)$ around an old synthetic center $x$.
\end{proposition}
\begin{proof}
The Rouch\'e disk contains exactly one zero. The rigorously indexed zero also lies there by the ball-containment assumption, so the two coincide. The triangle bound $|\gamma_k-x|\le r_1+|b_1-x|$ is only an error estimate and does not assert the indexed zero belongs to the smaller certified disk. Thus the enlarged-disk membership check does not justify the index identification.
\end{proof}
\begin{remark}[Validated arithmetic, not midpoint floats]\label{rem:interval}
All inequalities in \CGeqref{eq:oneineq} and \CGeqref{eq:samedisc} must be decided by strict outward-rounded ball bounds. Converting certified midpoints and radii to Python floating-point values for a \emph{decision} can discard the very margin being certified. Midpoints may be used to propose new exact rational or dyadic centers, but the final tests must be rerun using rigorous balls.
\end{remark}
\begin{corollary}[Certified real zeros in a finite set of discs]\label{cor:finite}
If $n$ distinct indices possess validated one-zero Rouch\'e discs and rigorously indexed Arb enclosures satisfy \CGeqref{eq:samedisc} for each index, then each of those $n$ indexed zeros is real and simple. This remains a finite validation statement; neither a uniform-in-height theorem nor RH follows.
\end{corollary}

\section{Hankel moments, block reality, and isolated close pairs}\label{sec:clusters}
A close pair can be certified without being separated by a synthetic clock, but a global proof must also cover every exceptional block. The two mechanisms below are independent of the arithmetic flow and involve the true completed function.
\begin{theorem}[Moment criterion for zeros enclosed by a contour]\label{thm:hankel}
Let $g$ be a real entire function, and let $\Omega$ be a bounded conjugation-invariant domain with zero-free boundary enclosing $n\ge1$ zeros, counted with multiplicity. For $b\in\R$ put
\begin{equation}\label{eq:moments}
 m_j=\frac1{2\pi i}\oint_{\partial\Omega}(z-b)^j\frac{g'(z)}{g(z)}\dd z,
 \qquad H_n=(m_{r+s})_{0\le r,s<n}.
\end{equation}
Then the real symmetric matrix $H_n$ satisfies
\begin{equation}\label{eq:hankelcriteria}
 H_n\succeq0\iff\text{every enclosed zero is real},\qquad
 H_n\succ0\iff\text{all enclosed zeros are real and simple}.
\end{equation}
When $n=2$, the enclosed zeros are
\begin{equation}\label{eq:tworoot}
 z_\pm=b+\frac{m_1\pm\sqrt{2m_2-m_1^2}}2.
\end{equation}
\end{theorem}
\begin{proof}
The residue theorem gives $m_j=\sum_{\nu=1}^{n}(z_\nu-b)^j$, including multiplicities. Since the domain is conjugation invariant, $m_j$ is real. For real coefficients $v=(v_0,\ldots,v_{n-1})$ write $P(w)=\sum_{j=0}^{n-1}v_jw^j$. Then
\begin{equation}\label{eq:hankelQF}
 v^{\mathsf T}H_nv=\sum_{\nu=1}^{n}P(z_\nu-b)^2.
\end{equation}
If the roots are real, the sum is nonnegative and is positive for every nonzero $P$ precisely when there are $n$ distinct real roots. If a conjugate nonreal pair is present, interpolate a polynomial with real coefficients which takes the values $i,-i$ on that pair and zero on all other distinct roots. This is possible with degree smaller than the number of distinct roots, at most $n$, using conjugate-symmetric Lagrange interpolation. The sum is then negative, proving failure of semidefiniteness. For two roots, the quadratic formula and the power sums give \CGeqref{eq:tworoot}.
\end{proof}
\begin{remark}\label{rem:hankel}
This is a contour version of the classical Hermite criterion. If $\Omega$ does not respect conjugation, positive definiteness need not imply reality. For RH, a family of positive-semidefinite matrices is sufficient only if it covers \emph{every} zero of $f(z)=\xi(1/2+iz)$; establishing such global positivity is an unproved RH-level condition, not a consequence of synthetic positivity in Proposition~\CGref{prop:loewner}.
\end{remark}

\begin{theorem}[Fourth-order localization under analytic isolation]\label{thm:isolated}
Let $t_*\in\R$, $0<h<R$, and suppose $g$ is real holomorphic on a neighborhood of $|z-t_*|\le R$ with
\begin{equation}\label{eq:factorH}
 g(z)=((z-t_*)^2-h^2)H(z),\qquad
 0<m_*\le|H(z)|\le M_*,\quad B=M_*/m_*,
\end{equation}
where $H$ is holomorphic and nonvanishing on the closed disk. If
\begin{equation}\label{eq:isolcond}
 \frac hR\le\frac1{1000B},
\end{equation}
then $g'$ has precisely one real zero $t_b$ in $[t_*-h/4,t_*+h/4]$, and
\begin{equation}\label{eq:tb}
 |t_b-t_*|\le1.01 B\frac{h^2}R,\qquad g(t_b)g''(t_b)<0.
\end{equation}
With $h_b=(-2g(t_b)/g''(t_b))^{1/2}$,
\begin{equation}\label{eq:hb}
 |h_b-h|\le10B^2\frac{h^3}{R^2},\qquad
 \frac{h_b^3}{|g''(t_b)|}\sup_{|z-t_b|\le2h_b}|g^{(5)}(z)|
 \le520B\left(\frac hR\right)^3.
\end{equation}
The fourth-order Taylor polynomial $T_4$ of $g$ around $t_b$ possesses two distinct real simple roots $y_\pm$ satisfying
\begin{equation}\label{eq:q4}
 |y_\pm-(t_*\pm h)|\le40B\frac{h^4}{R^3},\qquad
 \frac{|y_\pm-(t_*\pm h)|}{2h}\le20B\left(\frac hR\right)^3.
\end{equation}
\end{theorem}
\begin{proof}
Write $x=z-t_*$. Cauchy's estimates in the half-radius disk show
\begin{equation}\label{eq:cauchyH}
 |H^{(j)}(z)|\le j!M_*\left(\frac2R\right)^j
 \qquad(|z-t_*|\le R/2).
\end{equation}
Differentiating the factorization gives
\begin{equation}\label{eq:f2H}
 g'=2xH+(x^2-h^2)H',\quad
 g''=2H+4xH'+(x^2-h^2)H''.
\end{equation}
For real $|x|\le h/4$, write $g''=2H(1+\eps)$; then
$|\eps|\le Bh/R+4.25B h^2/R^2<1.01\cdot10^{-3}$. Hence $g''$ has a fixed sign there and $|g''|\ge1.99m_*$. Since $|g'(t_*)|\le2M_*h^2/R$ and the change of $g'$ across each half-interval exceeds $0.49m_*h$, there is one and only one root $t_b$. The mean-value theorem gives the first bound in \CGeqref{eq:tb}. With $\delta=t_b-t_*$, $g(t_b)=(\delta^2-h^2)H(t_b)$ and the fixed sign of $g''$ proves the sign claim. The identity
\begin{equation}\label{eq:hbsq}
 h_b^2=\frac{h^2-\delta^2}{1+\eps(t_b)}
\end{equation}
then yields $|h_b-h|\le10B^2h^3/R^2$.

By Leibniz' rule,
\begin{equation}\label{eq:f5}
 g^{(5)}=20H^{(3)}+10(z-t_*)H^{(4)}+((z-t_*)^2-h^2)H^{(5)}.
\end{equation}
The disk $|z-t_b|\le2h_b$ is contained in $|z-t_*|\le3h\le R/2$. Applying \CGeqref{eq:cauchyH} to the three terms gives $\sup|g^{(5)}|\le1000M_*/R^3$. Since $h_b\le1.01h$ and $|g''(t_b)|\ge1.99m_*$, the second inequality in \CGeqref{eq:hb} follows.

Finally choose $r=40Bh^4/R^3\ll h$ and the two disjoint disks $D_\pm=\{|z-(t_*\pm h)|\le r\}$. On their boundaries, the fourth-order Taylor remainder satisfies
\begin{equation}\label{eq:R5}
 |g-T_4|\le\frac{1000M_*}{R^3}\frac{(1.002h)^5}{120}\le8.5M_*\frac{h^5}{R^3},
\end{equation}
whereas the exact factorization gives
\begin{equation}\label{eq:comp4}
 |g(z)|\ge m_*r(2h-r)\ge79.6M_*\frac{h^5}{R^3}.
\end{equation}
Rouch\'e's theorem applied in the direction $|g-T_4|<|g|$ shows that $T_4$ has exactly one root in each $D_\pm$, as $g$ does. Since the polynomial has real coefficients and each disk is conjugation invariant, the roots are real and simple. The chosen radius gives \CGeqref{eq:q4}.
\end{proof}
\begin{remark}[What the isolation hypothesis does not say]\label{rem:isolated}
The theorem is conditional on the nonvanishing and two-sided bounds on $H$ \emph{throughout the specified complex disk}. A small real separation $2h$ alone does not imply $h/R\ll1$. Even arbitrarily accurate location of infinitely many real pairs would not exclude additional off-axis zeros elsewhere.
\end{remark}

\begin{proposition}[A limiting spacing law does not eliminate exceptional coincidences]\label{prop:portmanteau}
Let $s_k\ge0$ be normalized nonnegative consecutive ordinate gaps. If $N^{-1}\sum_{k\le N}\delta_{s_k}\Rightarrow\mu$ and $\mu(\{0\})=0$, then the set of indices with $s_k=0$ has density zero. Conversely, changing $s_k$ on a zero-density set, even inserting infinitely many exact zeros, does not alter the weak limit.
\end{proposition}
\begin{proof}
For every $\eps>0$, by the portmanteau theorem applied to the closed set $[0,\eps]$,
\begin{equation}\label{eq:portmanteau}
 \limsup_N\frac{\#\{k\le N:s_k=0\}}N\le\mu([0,\eps]).
\end{equation}
Let $\eps\downarrow0$. For the converse, bounded continuous test functions differ in their two empirical averages by at most $2\|\varphi\|_\infty$ times the exceptional density, which tends to zero.
\end{proof}
\begin{remark}[Exact local counterexample]\label{rem:quartic}
For $h>0$, $g_h(z)=(z^2-h^2)(z^2+4h^2)$ has two real roots $\pm h$ and two nonreal roots $\pm2ih$, although $g_h^{(5)}\equiv0$. Thus perfect fourth-order Taylor approximation at the real pair cannot by itself exclude a nonreal pair. The factor $H(z)=z^2+4h^2$ vanishes at $\pm2ih$, preventing the isolation hypothesis \CGeqref{eq:factorH} from holding on a disk enclosing all four zeros.
\end{remark}

\section{Reported experiments and independent arithmetic checks}\label{sec:checks}
The author's repository contains the original \texttt{rouche\_one\_zero.py}
validator (using \texttt{python-flint}; full path below) and, as of the
23 October 2026 commit \texttt{011fb1984}, the completed CSV, JSON, and log files
for two blocks of 300 consecutive indices each. The reported figures are
reproduced from those archived outputs, not from a new Arb execution. In
particular, the original program uses \texttt{acb.zeta\_zero(k)} as a
Turing-indexed comparator and does not read \texttt{zeross.bin} while
constructing the test centers.
\begin{center}\small
\begin{tabular}{@{}lrrrr@{}}\toprule
Indices & Tested & Local Rouch\'e pass & Test failed & Index mismatches\\\midrule
$1\,000\,055$--$1\,000\,354$ & 300 & 292 & 8 & 0\\
$1\,400\,737$--$1\,401\,036$ & 300 & 292 & 8 & 0\\\midrule
Total & 600 & 584 & 16 & 0\\\bottomrule
\end{tabular}
\end{center}
Of the 584 reported successful local tests, 346 used the original center
and 238 used Newton recentering. Reported mean radii about the original
prime-corrected centers were $0.0384809$ and $0.0429626$, with maxima
$0.1366163$ and $0.1694444$, for the first and second block respectively;
the second block recorded two overlapping consecutive \emph{reported}
error discs, and the first recorded none. The initial ten-index pilot
reported $10/10$ with mean radius $0.0509514$.

\textbf{Status of the certification.} The strict Taylor--Rouch\'e
inequality is evaluated in Arb balls and, where it succeeds, gives a
\emph{local existence-and-simplicity result} provided the holomorphic jet
and disc remainder enclosure are correct. Nevertheless, the original
index check has a material flaw after Newton recentering: it compares
the output of \texttt{acb.zeta\_zero(k)} with the original center $x$ and
the \emph{enlarged} error radius $r_1+|b_1-x|$, not with the actual
certified Rouch\'e disc $D(b_1,r_1)$. Membership in the larger disc does
not imply membership in the smaller. Furthermore, its ordinal-membership
and reality comparisons are performed through \texttt{float} midpoints
and radii, not strict ball containment. Consequently, the CSV status
\texttt{ok} in those 238 recentered cases is \emph{not yet a rigorous
certification of the advertised ordinal index}, even though the local
Rouch\'e inequality may succeed. The original-center tests likewise
require interval rather than \texttt{float} comparisons for a formal
end-to-end indexed certificate. The independent amended validator is
specified in Appendix~\CGref{app:onecodesource}, but has not been
executed here. No other inference about the untested indices follows.

The source records are at
\url{https://github.com/vcasana1/LawCasanya/tree/main/work/paper1/rouche_cert}
and their publication commit is
\url{https://github.com/vcasana1/LawCasanya/commit/011fb1984}.

\subsection{Two archived two-root examples}\label{subsec:two-checks}
The independent high-precision ordinary-arithmetic re-evaluation of the \emph{central values} of $f(b)$, $f'(b)$ and $f''(b)$ reproduces the published margins at the following heights. The tabulated $M_3$ bounds are \emph{inputs transcribed from the earlier interval program, not independently recertified by Arb/FLINT}.
\begin{center}\small
\begin{tabular}{@{}lrr@{}}\toprule
Quantity, after stated positive rescaling of $f$ & Height $1283$ & Height $1290$\\\midrule
$b$ & $1283.159048033$ & $1290.255802206$\\
$h$ & $0.164715498$ & $0.154422788$\\
Scaling multiplier & $10^{433}$ & $10^{435}$\\
$|f''(b)|$ (ordinary precision) & $52.034634808$ & $18.990581971$\\
Reported $M_3$ upper bound & $180.855125215$ & $63.282575959$\\
Small-circle left side (conditional) & $0.26309525$ & $0.07585725$\\
Small-circle right side & $0.30882290$ & $0.09906246$\\
Large-circle left side (conditional) & $0.45462859$ & $0.13108133$\\
Large-circle right side & $0.88235115$ & $0.28303560$\\
$hM_3/|f''(b)|$ (conditional) & $0.57249642$ & $0.51458517$\\\bottomrule
\end{tabular}
\end{center}
Both central checks satisfy the numerical inequalities if the previously reported $M_3$ bounds are correct. The corresponding exact complex-disk $M_3$ audit remains open because an independent Arb/FLINT execution was not available in the present environment. These numbers therefore are \emph{high-precision cross-checks}, not newly completed Arb certificates.

\subsection{Required reproducible audit protocol}\label{subsec:audit}
An independent audit of the two examples must compute \emph{all} of the following using outward-rounded balls: (i) validated $f(b)$, $f'(b)$ and $f''(b)$, (ii) a certified upper bound for $M_3$ on the \emph{whole complex disk} $|z-b|\le3h/2$, for instance from interval third-derivative jets covering the boundary and the maximum-modulus principle, and (iii) strict interval versions of \CGeqref{eq:I} and \CGeqref{eq:II}. It must retain a machine-readable log, commit the script with the exact output, and record the Arb and FLINT versions. A high-precision sample of finitely many boundary points is not sufficient to establish $M_3$.

For the one-zero experiment, the audit must additionally verify the precise \texttt{acb.zeta\_zero(k)} ordinal guarantee, use ball inclusion in the \emph{same certified disk}, check analytic branches, and distinguish radii about the final Rouch\'e center from error radii about the original prime-only center. The source of centers must be documented and must not use \texttt{zeross.bin}.

\section{Global research barriers}\label{sec:status}
The local two-zero and one-zero theorems are classical applications of Rouch\'e made effective with arithmetically proposed centers. In particular, neither certification of arbitrarily many real zeros nor a success fraction approaching one implies that all zeros are on the critical line. A hypothetical off-line pair would produce two zeros sharing the same ordinate and could constitute an exceptional set invisible to averaged displacement or limiting spacing laws.

To complete the direct ordinate route, one must derive \CGeqref{eq:adjacent} or \CGeqref{eq:directional} \emph{for every index}, from new arithmetic estimates not assuming the distinctness of the true ordinates. To complete the Rouch\'e route, one must obtain a theorem of universal reality plus exhaustive coverage of all potential zeros, not merely the real roots detected by a locally successful test. Reducing coverage to the full classical Turing zero count without an independent global inequality is a known RH-level obstruction. The positive Loewner kernel \CGeqref{eq:Loewner} is an unconditional statement about the synthetic product and has no established positive transfer to $\xi$.

The spectral-pair isolation program remains a separate legitimate problem: relative isolation $h/R\ll1$ and a controlled nonzero auxiliary factor are more informative than merely the smallness of an absolute gap $h$. No universal such isolation bound is derived here, and statistical GUE repulsion does not exclude exceptional multiple or off-critical zeros. The stationary rational-frequency expansion \CGeqref{eq:fourier} is exact and prime-only, but uniform cancellation across growing cutoffs remains unproved. We do not interpolate synthetic $\zeta_c$ to true $\zeta$ as a purported RH proof; the endpoint transfer problem is still uncontrolled.

\begin{thebibliography}{99}

\bibitem{DelvesLyness1967} L.~M. Delves and J.~N. Lyness, \emph{A numerical method for locating the zeros of an analytic function}, Math. Comp. \textbf{21} (1967), 543--560. \href{https://doi.org/10.1090/S0025-5718-1967-0228165-4}{doi:10.1090/S0025-5718-1967-0228165-4}.

\bibitem{Edwards1974} H.~M. Edwards, \emph{Riemann's Zeta Function}, Academic Press, 1974.

\bibitem{Gonek2012} S.~M. Gonek, \emph{Finite Euler products and the Riemann hypothesis}, Trans. Amer. Math. Soc. \textbf{364} (2012), 2157--2191. \href{https://doi.org/10.1090/S0002-9947-2011-05546-7}{doi:10.1090/S0002-9947-2011-05546-7}.

\bibitem{KravanjaVanBarel2000} P.~Kravanja and M.~Van Barel, \emph{Computing the Zeros of Analytic Functions}, Springer, 2000. \href{https://doi.org/10.1007/BFb0103927}{doi:10.1007/BFb0103927}.

\bibitem{KreinNaimark1981} M.~G. Krein and M.~A. Naimark, \emph{The method of symmetric and Hermitian forms in the theory of the separation of the roots of algebraic equations}, Linear and Multilinear Algebra \textbf{10} (1981), 265--308. \href{https://doi.org/10.1080/03081088108817420}{doi:10.1080/03081088108817420}.

\bibitem{Mangoldt1905} H.~von Mangoldt, \emph{Zur Verteilung der Nullstellen der Riemannschen Funktion $\xi(t)$}, Math. Ann. \textbf{60} (1905), 1--19. \href{https://doi.org/10.1007/BF01447494}{doi:10.1007/BF01447494}.

\bibitem{MontgomeryVaughan2006} H.~L. Montgomery and R.~C. Vaughan, \emph{Multiplicative Number Theory I: Classical Theory}, Cambridge University Press, 2006. \href{https://doi.org/10.1017/CBO9780511618314}{doi:10.1017/CBO9780511618314}.

\bibitem{Titchmarsh1986} E.~C. Titchmarsh, \emph{The Theory of the Riemann Zeta-Function}, 2nd ed., revised by D.~R. Heath-Brown, Clarendon Press, 1986.

\bibitem{Turing1953} A.~M. Turing, \emph{Some calculations of the Riemann zeta-function}, Proc. London Math. Soc. (3) \textbf{3} (1953), 99--117. \href{https://doi.org/10.1112/plms/s3-3.1.99}{doi:10.1112/plms/s3-3.1.99}.

\end{thebibliography}

\appendix
\section{A minimal independent numerical cross-check}\label{app:precheck}
The following executable Python program reproduces the table's central derivative checks. It uses \texttt{mpmath} high precision; by design it does \emph{not} claim rigorous outward-rounded bounds for $M_3$.
\begin{lstlisting}[language=Python]
import mpmath as mp
mp.mp.dps = 65
cases = [
 ('1283.159048033','0.164715498','1e433','180.855125215'),
 ('1290.255802206','0.154422788','1e435','63.282575959')]
def F(z):
    s = mp.mpf('0.5') + 1j*z
    return s*(s-1)/2 * mp.pi**(-s/2) * mp.gamma(s/2)*mp.zeta(s)
for bs,hs,scale,M3s in cases:
    b,h,t,M3=map(mp.mpf,(bs,hs,scale,M3s))
    a=t*F(b); a1=t*mp.diff(F,b,1); a2=t*mp.diff(F,b,2)
    err=abs(a+a2*h*h/2)
    sl=err+5*h*abs(a1)/4+mp.mpf(125)/384*M3*h**3
    sr=7*abs(a2)*h*h/32
    bl=err+3*h*abs(a1)/2+mp.mpf(9)/16*M3*h**3
    br=5*abs(a2)*h*h/8
    print(bs, 'small',sl,sr,'outer',bl,br)
\end{lstlisting}

\section{Algorithmic validation requirements for Arb/FLINT}\label{app:arb}
For future rigorous validation the one-zero calculation should follow this logic. The following is a \emph{procedure specification}, not an assertion that a complete independent Arb audit has been executed:
\begin{enumerate}[leftmargin=*]
\item Read the synthetic center $x_k$ as the exact dyadic rational represented by its binary64 input; compute Taylor coefficient balls of locally holomorphic $Z$ on that center.
\item Choose an exactly represented center $b$ (either $x_k$ or a Newton candidate) and exactly represented radius $r$; enclose each $a_j$ and $M_K$ using Arb balls, and certify the strict inequality \CGeqref{eq:oneineq}.
\item Obtain the library's certified $k$th ordinate enclosure $\mathcal B_k$ and verify its documented ordinal meaning; test \CGeqref{eq:samedisc} using directed ball endpoints, \emph{not} float midpoints.
\item If the radius is reported relative to $x_k$, return both the actual Rouch\'e disk $(b,r)$ and the derived bound $r+|b-x_k|$. Never use the latter radius to assert membership in the former disk.
\item Record the precision, package version, failure reason, exact inputs, resulting ball endpoints, and proof margins in the output. Unverified, overlapping, and failed disks must be retained as such.
\end{enumerate}
For a two-zero certificate, enclose $f'''$ uniformly on the boundary circle $|z-b|=3h/2$ using finite interval \emph{arcs}, not sample nodes; maximum modulus then bounds the whole disk. The two inequalities \CGeqref{eq:I} and \CGeqref{eq:II} must be strict in ball arithmetic. An implementation template, recorded separately with this manuscript, is intended to support this reproducibility step.

\section{Independent Arb/FLINT pair-audit source (not executed here)}\label{app:arbsource}
The following program is provided for reproduction. Its local syntax was checked, but
its numerical Arb execution has not been performed in this environment. A valid
certificate requires a successful run with retained interval output.
\begin{lstlisting}[language=Python]
"""Independent Arb/FLINT audit template for the two Casanya pair discs.

Python >=3.10, python-flint with acb_series.gamma() and acb_series.zeta().
No genuine ordinates, zeross.bin or reference zero list are used.

The certificate is valid only if this program completes with a strict PASS for
both cases and its interval enclosures are retained. NOT EXECUTED in the
ChatGPT container on 2026-10-23: python-flint was unavailable there.

Run: python audit_pairs_arb.py --case 1283 --arcs 4096
     python audit_pairs_arb.py --case 1290 --arcs 4096

Algorithm: finite interval arcs cover the *entire boundary circle* of the
radius 3*h/2. Each interval third derivative is enclosed via the coefficient
of w**3 in a Taylor series of F(z+w); maximum modulus covers the disk.
The Taylor coefficients of the real entire function Xi(1/2+i*z) are enclosed
by ball evaluation at exact decimal input b and a positive scaling constant.
No assumed global bounding constant or FFT tail estimate is used.
"""
import argparse
import json
import platform
import time
from flint import acb, arb, acb_series, ctx

CASES = {
    1283: {"b": "1283.159048033", "h": "0.164715498", "scale": "1e433"},
    1290: {"b": "1290.255802206", "h": "0.154422788", "scale": "1e435"},
}

def scaled_xi_jet(zball, scale):
    """Taylor coefficients through 3, each a complex ball.

    Parameter zball may be an acb box; if so the coefficient j encloses
    F^(j)(z)/j! for every z in that box.
    """
    x = acb_series([zball, acb(1)], prec=4)
    s = acb("0.5") + acb(0, 1) * x
    gamma_factor = (s / 2).gamma()
    completed = (s * (s - 1) / 2) * ((-s / 2) * arb.pi().log()).exp()
    z = acb(scale) * completed * gamma_factor * s.zeta()
    coeff = z.coeffs()
    if len(coeff) < 4:
        coeff.extend([acb(0)] * (4 - len(coeff)))
    return coeff[:4]

def interval_arc(j, n):
    """Outward-rounded angle interval covering [2*pi*j/n,2*pi*(j+1)/n]."""
    pi = arb.pi()
    mid = pi * arb(2 * j + 1) / n
    half_width = (pi / n).upper()
    radius = half_width + mid.rad()
    return arb(mid.mid(), radius)

def audit(case, arcs, bits):
    ctx.prec = bits
    ctx.cap = 16
    cfg = CASES[case]
    b, h, scale = arb(cfg["b"]), arb(cfg["h"]), arb(cfg["scale"])
    radius = arb("1.5") * h
    start = time.time()
    central = scaled_xi_jet(acb(b), scale)
    f0, f1, f2 = central[0], central[1], 2 * central[2]
    # Strict realness of Xi along R is known analytically. The absolute
    # complex ball bounds are safe even if a tiny imaginary radius survives.
    f2abs = abs(f2)
    f2low = f2abs.lower()
    if not (f2low > 0):
        raise ArithmeticError("f'' lower bound does not exclude zero")
    bound = arb(0)
    for j in range(arcs):
        theta = interval_arc(j, arcs)
        zarc = acb(b + radius * theta.cos(), radius * theta.sin())
        j3 = scaled_xi_jet(zarc, scale)[3]
        upper = abs(6 * j3).upper()
        if upper > bound:
            bound = upper
        if j % max(1, arcs // 16) == 0:
            print("case", case, "arc", j, "/", arcs, "M3upper", bound, flush=True)
    # Holomorphy of xi implies |f'''(z)| <= max_boundary |f'''|.
    M3 = bound
    residual = abs(f0 + f2 * h * h / 2)
    H = h.upper()
    # Bounds computed monotonically with directed interval endpoints.
    small_left = residual.upper() + (arb(5) * H / 4) * abs(f1).upper() + arb(125) * M3 * H**3 / 384
    small_right = arb(7) * f2low * h.lower()**2 / 32
    outer_left = residual.upper() + (arb(3) * H / 2) * abs(f1).upper() + arb(9) * M3 * H**3 / 16
    outer_right = arb(5) * f2low * h.lower()**2 / 8
    eta = H * M3 / f2low
    small_ok = bool(small_left.upper() < small_right.lower())
    outer_ok = bool(outer_left.upper() < outer_right.lower())
    status = "PASS" if small_ok and outer_ok else "INCONCLUSIVE"
    obj = dict(
        status=status, case=case, input=cfg, arcs=arcs, bits=bits,
        M3_upper=str(M3), f2_abs_lower=str(f2low),
        small_left_upper=str(small_left.upper()), small_right_lower=str(small_right.lower()),
        outer_left_upper=str(outer_left.upper()), outer_right_lower=str(outer_right.lower()),
        eta_upper=str(eta.upper()),
        central_coefficients=[str(x) for x in central],
        elapsed_s=round(time.time() - start, 2),
        python=platform.python_version(),
    )
    print(json.dumps(obj, indent=2), flush=True)
    with open(f"pair_{case}_arb_audit.json", "w") as o:
        json.dump(obj, o, indent=2)
    if status != "PASS":
        raise RuntimeError("No certificate: increase precision/arcs or improve the interval bounds")

if __name__ == "__main__":
    ap = argparse.ArgumentParser()
    ap.add_argument("--case", type=int, choices=[1283, 1290], required=True)
    ap.add_argument("--arcs", type=int, default=4096)
    ap.add_argument("--bits", type=int, default=256)
    args = ap.parse_args()
    audit(args.case, args.arcs, args.bits)
\end{lstlisting}

\section{Index-corrected one-zero Arb validator (not executed here)}\label{app:onecodesource}
The following implementation amends the same-disc ordinal matching test of
Proposition~\CGref{prop:index}. The file is also supplied separately. The
code is syntax checked, not executed here: the required FLINT library was
unavailable in this environment.
\begin{lstlisting}[language=Python]
"""Rouche + indexed Turing/Arb one-zero verifier (corrected validation logic).

Requires python-flint, numpy; no zeross.bin or true zeros in construction.
Usage: python rouche_one_zero_v4.py CENTERS.bin K0 COUNT OUT_PREFIX

This program has been syntax checked but *not* run in the ChatGPT environment
because python-flint is unavailable. All comparisons use Arb ball bounds.
It uses acb.zeta_zero(k) only after a purely local Rouche disk is validated.
"""
import sys
import json
import time
import numpy as np
from flint import acb, arb, acb_series, ctx

ctx.prec = 160
ctx.cap = 16
PI = arb.pi()
K = 7
RMAX = arb("0.2")
TRIES = ("1.5", "2", "3", "5")

def jet(b):
    """Hardy Z Taylor coefficients of degrees 0 through K at a ball b."""
    x = acb_series([acb(b), acb(1)], prec=K + 2)
    s = acb_series([acb("0.5") + acb(0, 1) * acb(b), acb(0, 1)], prec=K+2)
    u = acb_series([acb("0.25") + acb(0, 1) * acb(b) / 2, acb(0, 1) / 2], prec=K+2)
    uc = acb_series([acb("0.25") - acb(0, 1) * acb(b) / 2, acb(0, -1) / 2], prec=K+2)
    theta = (u.lgamma() - uc.lgamma()) / acb(0, 2) - x * (PI.log() / 2)
    values = ((theta * acb(0, 1)).exp() * s.zeta()).coeffs()
    if len(values) < K + 1:
        values.extend([acb(0)] * (K + 1 - len(values)))
    return values

def local_disc(center):
    """Return a disk certified by Rouche, or None.

    CENTER is an exact dyadic/input Arb real ball; outward interval bounds
    are used for all decisive comparisons.
    """
    z = jet(center)
    a, c = abs(z[0]), abs(z[1])
    if not (c.lower() > 0):
        return None
    r0 = a / c
    for m in TRIES:
        radius = (arb(m) * r0).upper() + arb("1e-12")
        if not (radius.upper() < RMAX.lower()):
            continue
        # This complex rectangle contains the whole circle |z-center|<=radius.
        box = acb(arb(center, radius), arb(0, radius))
        zK = abs(jet(box)[K]).upper()
        high_r = radius.upper()
        remainder = sum((abs(z[j]).upper() * high_r**j for j in range(2, K)), arb(0))
        remainder += zK * high_r**K
        lhs = a.upper() + remainder.upper()
        rhs = c.lower() * radius.lower()
        if bool(lhs < rhs):
            return dict(center=center, radius=radius, lhs=lhs, rhs=rhs,
                        multiplier=m, f0_abs=a, f1_abs=c, rem=remainder)
    return None

def certify_index(k, certified):
    """Verify inclusion of the Arb Turing indexed zero in THE certified disk."""
    gamma_ball = acb.zeta_zero(k)
    center = certified["center"]
    radius = certified["radius"]
    zeta_s_center = acb(arb("0.5"), center)
    upper_distance = abs(gamma_ball - zeta_s_center).upper()
    return bool(upper_distance < radius.lower()), gamma_ball, upper_distance

def run(path, k0, count, prefix):
    xs = np.fromfile(path, dtype="<f8")
    records = []
    start = time.time()
    for k in range(k0, k0 + count):
        x = float(xs[k-1])
        xb = arb(x)  # the binary64 input is only a proposed, prime-only center
        trial = local_disc(xb)
        method = "original"
        if trial is None:
            v = jet(xb)
            if not (abs(v[1]).lower() > 0):
                records.append({"k":k,"status":"fail_derivative"})
                continue
            # Proposal only: round the Newton correction, then fully recertify.
            correction = (v[0]/v[1]).real.mid()
            bnew = arb(float(x - float(correction)))
            trial = local_disc(bnew)
            method = "newton"
        if trial is None:
            records.append({"k":k,"status":"fail_rouche"})
            continue
        ok, ball, distance = certify_index(k, trial)
        effective_bound = trial["radius"].upper() + abs(trial["center"]-xb).upper()
        records.append(dict(
            k=k, input_center=repr(x), certified_center=str(trial["center"]),
            actual_radius=str(trial["radius"]), bound_about_input=str(effective_bound),
            rouche_margin=str(trial["rhs"].lower() - trial["lhs"].upper()),
            ball_k=str(ball), ball_to_cert_center_upper=str(distance),
            status="indexed_certified" if ok else "index_not_proved", method=method))
        if (k-k0+1) % 10 == 0:
            print("processed", k-k0+1, "of", count, flush=True)
    summary = dict(input=path,k0=k0,count=count,
        indexed_certified=sum(r["status"] == "indexed_certified" for r in records),
        fail_rouche=sum(r["status"] == "fail_rouche" for r in records),
        index_not_proved=sum(r["status"] == "index_not_proved" for r in records),
        time_s=round(time.time()-start,2))
    with open(prefix+".json", "w") as out:
        json.dump(dict(summary=summary, records=records), out, indent=2)
    print(json.dumps(summary,indent=2))

if __name__ == "__main__":
    if len(sys.argv)!=5:
        raise SystemExit("Usage: rouche_one_zero_v4.py CENTERS.bin K0 COUNT OUT_PREFIX")
    run(sys.argv[1],int(sys.argv[2]),int(sys.argv[3]),sys.argv[4])
\end{lstlisting}

\section{Explicit limitations and open questions}\label{app:limitations}
\begingroup\small
\begin{enumerate}[leftmargin=*,nosep]
\item The original local experiments at heights $1283$ and $1290$ used an earlier interval engine; independently certified Arb third-derivative bounds are pending.
\item The published Arb run contains complete CSV/JSON/log summaries of two 300-index blocks: 292/300 local tests marked successful in each (584/600 overall); 16 attempts failed the chosen Rouch\'e test. This is a finite experiment, not a universal theorem.
\item The original Python prototype has an index-matching error after Newton recentering and uses floating-point membership comparisons; 238 of its 584 successful flags involve Newton recentering. The local Rouch\'e test and ordinal identification must be distinguished. Strict Arb containment checks in the \emph{actual} Rouch\'e discs remain to be rerun before the 584 flags may be reported as rigorously indexed certificates.
\item A rigorously indexed zeta zero returned by Arb already uses evaluation of the genuine zeta function and Turing's method. This is allowed in \emph{validation}, but does not make the prime-only construction itself an identification theorem.
\item A local two-zero certificate classifies roots in its certified disk. It cannot exclude nonreal roots in other parts of the strip.
\item The order-preserving prime flow may fail to locate a particular ordinate within its own cell. It controls synthetic order, not RH.
\item The choice $a=1/\sqrt2$ is not known to minimize individual errors or their worst-case normalized gap ratios.
\item The observed smallness of $|\gamma_k-c_k|$ for many indices does not imply its global bound at all indices; an isolated violation remains possible.
\item Statistical spacing limits, even if proved in full, do not exclude zero-density sets of repeated ordinates.
\item Exact synthetic functional equations and rational-frequency formulas do not identify the synthetic spectrum with the true one.
\item No unconditional positivity estimate transferring the Loewner kernel from $\Xi_c$ to $\xi$ is available.
\item Neither the local certificate nor the arithmetic flow constitutes a proof of RH or of universal simplicity.
\end{enumerate}
\endgroup


\end{document}
